\scanpage{131}

\[
\deleted{\unsure{Z}}\,G_0=\mathrm{Aut}(S,F_0)=\mathrm{Aut}(S,A_0)=\mathrm{Aut}\,P_0\tag*{(5)}
\]
(où $P_0=(S,A_0,F_0,\text{ relations d'incidences}\ldots)$)
\[
\widehat Z=\mathrm{Centr}_{\unsure{\mathfrak S|\mathcal S}_S}(G_0)\tag*{(6)}
\]
% slanted
\margin{-- \unsure{laisser|je laisse} tomber le (\unsure{2|3|5}) finalement !}

\deleted{\unsure{(7)}} et soit
\[
\begin{aligned}
Z^{\unsure{*|\sharp}}=\bigl\{\lambda\in\widehat Z\ \big|\ &\exists f\in\mathfrak S_S\ \text{t.q. }\forall\{s,t\}\inS,\ \text{on a}\\
&\{s,t\}\in A_0\Longleftrightarrow\{f^{-1}s,f^{-1}\lambda t\}\in A_0\bigr\}.
\end{aligned}\tag*{(7)}
\]

\deleted{i.e.} On peut dire aussi que $\lambda\in\widehat Z$ est dans $Z^*$ ssi il satisfait les deux conditions :
\[
\left\{\begin{aligned}
\mathrm a)\quad &\forall(s,t)\in S,\qquad\{s,t\}\in A_0\Longleftrightarrow\{\lambda^{-1}s,\lambda t\}\in A_0\\
&[\,\text{d'où un } A_\lambda\subset\unsure{\mathfrak P|\mathcal P|P}_2(S),\ \text{défini de façon unique}\\
&\text{par la condition que }\{s,t\}\in A_0\Longleftrightarrow\{s,\lambda t\}\in A_\lambda\,]\\
\mathrm b)\quad &\text{le graphe }\Gamma_\lambda=(S,A_\lambda)\text{ est isomorphe}\\
&\text{au graphe }\Gamma_0=(S,A_0)
\end{aligned}\right.\tag*{(8)}
\]
\margin{Cela implique que $\Gamma_\lambda$ est associé à une structure \unsure{triangulée|triangul.|triangulaire} $F_\lambda$ sur $S$.}

On suppose

(9)\qquad (To 3)\qquad L'ensemble $Z^*\subset\widehat Z\subset\deleted{\unsure{G_0}}\mathfrak S_S$ est formé de \deleted{\unsure{t.tes|les}} permutations qui commutent deux à deux, \deleted{\unsure{ou|on}}

On voit (en distinguant le cas où tt sommet est déterminé par son étoile, et celui où il n'en est rien i.e. où $P_0$ est une double pyramide) que la condition (To 3)

% archive sheet number: 130
